\section{Introduction}
\label{sec:introduction}

The rapid advances in large language models for code generation have raised a practical question for the high-performance computing (HPC) community: can LLMs reliably translate parallel code across GPU programming APIs? Such translation has clear practical value. Scientific applications written in CUDA are increasingly needed to support AMD and Intel accelerators via OpenMP, OpenCL, HIP, or SYCL, and the manual cost of porting non-trivial kernels is substantial. 
% Experienced GPU programmers spend days to weeks ensuring thread-index arithmetic is correct, memory access patterns are coalesced, and synchronization semantics are preserved across API boundaries.

Yet parallel code translation is qualitatively different from sequential code tasks. Thread synchronization, shared memory management, and API-specific idioms (e.g., CUDA's \texttt{threadIdx}/\texttt{blockIdx} vs.\ OpenMP's \texttt{\#pragma omp parallel for}) require structural transformation, not surface-level renaming. Existing benchmarks do not address this. Sequential benchmarks such as HumanEval~\cite{HumanEval2021} and SWE-bench~\cite{SWEbench2024} contain no parallelism. Parallel code generation benchmarks, such as ParEval~\cite{ParEval2024}, evaluate synthesis from natural language descriptions---a different focus from translating existing parallel structure across APIs. Repository-level translation work (ParEval-Repo~\cite{ParEvalRepo2025}) finds 0\% pass@$k$ on applications larger than 133~SLoC, but the binding constraint is build-system generation, not parallel logic translation. To the best of our knowledge, prior frameworks have not explicitly distinguished genuine parallel reasoning from possible reliance on memorized translations of widely circulated benchmark code.

To address these gaps, we present \textbf{ParBench}, a kernel-centric benchmark framework for parallel code translation with a build-run-verify evaluation pipeline. ParBench isolates kernel translation from build-system generation by providing the surrounding build infrastructure as fixed context, directly targeting the failure mode identified by ParEval-Repo. Based on our systematic survey of 35 open-source HPC repositories, \textbf{ParBench} includes five benchmark suites (Rodinia, XSBench, RSBench, mixbench, and HeCBench), spanning six core translation directions over three API pairs (CUDA, OpenMP, and OpenCL), with two additional OpenMP-target directions used as case studies. 

To further evaluate whether LLMs reason about parallel structure rather than rely on surface similarity, an AST-driven augmentation engine applies six semantics-preserving transformations at five pre-defined intensity levels (L0--L4).

Specifically, ParBench is designed across following key dimensions:

\textbf{Multi-file coordination.} Beyond isolating kernel translation from build-system generation, ParBench evaluates multi-file coordination: 25\% of specifications require translating two or more source files. This reveals an independent difficulty dimension --- single-file translations achieve 51.3\% pass rate versus 22.2\% for multi-file translations ($\chi^2 = 82.73$, $p = 0.005$) --- demonstrating that file-boundary coordination is a substantial challenge distinct from within-file API mapping.

\textbf{Training data contamination --- the invisible gap.} A fourth gap cuts across all of the above: benchmark codes used in prior work are likely present in LLM training data, making it impossible to distinguish genuine parallel reasoning from memorized translations. Rodinia~\cite{Rodinia2009}, XSBench, and similar proxy applications have been republished across hundreds of GitHub repositories. ParBench addresses this directly through an AST-driven augmentation engine that produces semantically-equivalent but syntactically novel code variants, testing whether models reason about parallel semantics or simply pattern-match from training data (Section~\ref{sec:contamination}, Section~\ref{sec:augmentation-engine}). Empirical results confirm the practical importance of this concern (Section~\ref{sec:augmentation-robustness}).

\textbf{The gap in benchmark selection rationale.} A final dimension on which prior work is incomplete is the \emph{why} of benchmark selection: Which parallel APIs matter the most? Which kernels are representative? Existing frameworks do not answer these questions systematically. ParBench's selection is grounded in a comprehensive empirical survey of 35 open-source HPC repositories covering all the major parallel programming models. That survey identified 472 CUDA-OpenMP kernel pairs across just 6 repositories --- the largest available translation dataset in the ecosystem, and the practical bottleneck for real-world GPU-to-CPU portability work (Figure~\ref{fig:repo-vs-kernel} illustrates how repository-level counting understates kernel-level translation opportunities by up to two orders of magnitude). The survey further identified which benchmark suites provide the same kernel implemented across multiple APIs (Rodinia, HeCBench, XSBench, RSBench, mixbench), which have automatable build/run/verify pipelines, and which have self-checking output patterns. The choice of CUDA-to-OpenMP as the primary translation direction, and of Rodinia as the primary evaluation substrate, follows directly from this survey---not from convenience.

Together, these gaps define the problem that ParBench is designed to solve: kernel-level evaluation granularity, build-infrastructure isolation, multi-file coordination assessment, training-data robustness through augmentation, and survey-grounded benchmark selection.

This paper makes three contributions:

\begin{enumerate}

\item \textbf{Parallel code translation benchmark:} We introduce ParBench, a benchmark for LLM-based parallel code translation that isolates translation from repository-level build-system generation by fixing the surrounding infrastructure as context. This design enables direct build-run-verify evaluation of translation outcomes across parallel APIs.

% \item \textbf{ParBench framework}---We introduce ParBench, the first kernel-centric benchmark framework for LLM-based parallel code translation. It provides 96 executable specifications across five benchmark suites and six core translation directions over CUDA, OpenMP, and OpenCL, with two additional OpenMP-target directions as case studies. By fixing the surrounding build infrastructure as context, ParBench isolates kernel translation from repository-level build-system generation and enables direct evaluation of parallel translation capability at the kernel level.
\item \textbf{AST-driven augmentation engine:} We develop an AST-driven augmentation engine with six semantics-preserving source-level transformations applied at five augmentation levels. This engine enables controlled robustness evaluation under structural variation, making it possible to test whether LLMs preserve parallel API semantics beyond surface similarity while maintaining executability and functional correctness.


\item \textbf{Empirical characterization}---We provide the first systematic characterization of current LLM performance on parallel code translation under deterministic decoding and pass@$k$ sampling. Evaluation with ParBench demonstrates non-trivial translation capability under kernel isolation. It further shows that build-stage adaptation is the main bottleneck, performance is highly direction-dependent but broadly stable under augmentation, and current LLM capability on this task cannot be substantially improved through sampling alone.
\end{enumerate}



